\documentclass[10pt,a4paper]{amsart}
\usepackage[margin=19mm]{geometry}
\usepackage{amsmath,amssymb,mathtools}
\usepackage[scr=rsfs,cal=euler]{mathalfa}
\usepackage{xfrac}
\usepackage[unicode,hidelinks,bookmarks=false]{hyperref}
\newcommand{\ie}{i.e.\ }
\newcommand{\cf}{cf.\ }
\newcommand{\resp}{respectively\ }
\newcommand{\Subsection}{\subsection}
\providecommand{\nobreakdash}{\nobreak\hbox{-}}
\setlength{\emergencystretch}{2em}
\multlinegap0pt
\usepackage{amssymb}
\usepackage{cancel}
\usepackage{enumerate}
\newtheorem{thm}{Theorem}
\newtheorem{lem}{Lemma}
\newcommand{\D}{\partial}
\newcommand{\R}{\mathbb{R}}
\providecommand{\dv}{\,\mathrm{div}\,}
\title{Global existence for a 3D Tropical Climate Model with damping and small initial data in $\dot H^{\frac{1}{2}}(\R^3)$}
\author{Diego Berti, Luca Bisconti, Davide Catania}
\date{Source: arXiv:2305.07964v1 (CC BY 4.0)}
\begin{document}
\maketitle

\noindent\emph{Source and licence.} Global existence for a 3D Tropical Climate Model with damping and small initial data in $\dot H^{\frac{1}{2}}(\R^3)$. Version arXiv:2305.07964v1 (CC BY 4.0), distributed by arXiv under the Creative Commons Attribution 4.0 International licence, http://creativecommons.org/licenses/by/4.0/, as recorded in the arXiv licence metadata for this exact version and shown on the abstract page https://arxiv.org/abs/2305.07964v1. Full licence notice: \texttt{licenses/arxiv-2305.07964-CC-BY-4.0.txt}.\par\smallskip

\noindent\textit{Editorial scope.} Counted unit: the article's lem lem:monotone (source0.tex lines 1414--1448). The article's complete original proof of it is delivered with it (source0.tex lines 1477--1518, one \texttt{proof} environment of 42 lines, which the article places away from the statement). No mathematics is written, repaired or re-derived: the statement and the proof are the article's own lines, in the article's own notation, and no retained line is substituted or re-flowed.  Retained uncounted as the source-local context the proof needs: lines 129--137; lines 348--370; lines 1065--1069; lines 1399--1408.  Nothing was omitted from any retained span, so the package carries no omission record.  No retained line contains a citation, so the wrapper carries no bibliography and no bibliography entry.  No retained line contains a figure environment, an image-inclusion command or drawing code, so no image is delivered and the package declares no asset dependency.  Editorial formatting only, in addition: an AMS \texttt{amsart} preamble that restates the article's own declarations and macros used by the retained lines, namely the environments \texttt{thm}, \texttt{lem} on separate counters, together with the standard packages those lines need.  The retained environments print in this document as Theorem 1, Lemma 1 in order of appearance, whereas the article prints the same items with the numbering of its own full document.  The complete native source of the article as distributed in the arXiv e-print for this exact version is bundled unchanged as \texttt{sources/141449/source0.tex}.
\begin{equation}\label{TCM-gen}
  \begin{aligned}
    &\D_t u+ (u\cdot \nabla) u -\nu\Delta u + \sigma_1 |u|^{\alpha -1}u +\nabla \pi+\dv  (v \otimes v)=0,\\
    &\D_t v+ (u \cdot \nabla) v -\eta\Delta v + \sigma_2 |v|^{\beta -1}v + (v\cdot \nabla) u  +\nabla\theta =0,\\
    & \D_t \theta + (u \cdot \nabla) \theta  -\mu \Delta \theta  + \dv  v =0,\\
    &\dv  u = 0,\\
    & u (x, 0) =u_0,\,\, v (x, 0)=v_0,\,\, \theta (x, 0)=\theta_0,
  \end{aligned}
\end{equation}

    \begin{thm}[Blow-up criterion] \label{th:bu criterion}  Assume 
      $5/2\leq \alpha, \beta < 4$, and
      $(u_0, v_0, \theta_0) \in H^2(\R^3)^3$, with $\dv u_0 = 0$. Let $(u,v,\theta)(t)$ be the local solution of \eqref{TCM-gen}. Define
  \begin{equation}
  \label{e:Pi}
  \Pi(T) \doteq 
  \int_0^T C \big(1
      +\|u(t)\|_{{}_{\textrm{BMO}}}^{8} + \|v(t)\|_{{}_{\textrm{BMO}}}^{8} +
      \|\theta(t)\|_{{}_{\textrm{BMO}}}^2 +
      \|u(t)\|^{\alpha+1}_{\alpha+1}+\|v(t)\|^{\beta+1}_{\beta+1}
      \big)dt.
  \end{equation}
      
      Then , for $0< T^\ast < +\infty$, we have that  
  %
  \begin{equation*}
  \|\Delta (u, v, \theta) (T)\| < +\infty, \ \mbox{ for every } \ 0<T< T^\ast \ \mbox{ and } \ \limsup_{T \to {T^\ast}} \|\Delta (u, v, \theta) (T)\|=+\infty,
  \end{equation*}
  if and only if 
  \begin{equation*}
  \Pi(T) < +\infty \ \mbox{ for every } \ 0<T<T^\ast \ \mbox{ and } \ \Pi(T^\ast) =+\infty.
  \end{equation*}
  \end{thm}

\begin{equation} \label{stima-u-v-H1-chiusa}
  \frac{d}{dt} \|\nabla (u,v, \theta)(t)\|^2 + \Big(2\min\{\nu, \eta, \nu\} -
  C\big(\varepsilon + \|\Lambda^{\frac{1}{2}}
  (u,v,\theta)\|^2\big)\Big)\|\Delta (u, v, \theta)\|^2 \leq 0.
\end{equation}

\begin{equation} \label{e:key2}
  \begin{aligned}
    \frac{1}{2} \frac{d}{dt}  \|\Lambda^{\frac{1}{2}} & (u, v,
    \theta) (t) \|^2  +  \Phi(u,\alpha)+ \tilde \Phi (v,\beta)
    \\
    & + \big(\ell - C \|\Lambda^{\frac{1}{2}} (u, v, \theta) (t) \| -
    C\,\mathcal R(u,\alpha) -C\, \tilde{\mathcal R}(v,\beta)\big) \|\Lambda^{\frac{3}{2}} (u, v,
    \theta)\|^2 \leq 0.
  \end{aligned}
\end{equation}

\begin{lem}[Monotonicity of $\dot H^1$- and $\dot H^{\frac12}$-norms]\label{lem:monotone}
  Assume $5/2 \le \alpha, \beta <4$. Let $(u_0, v_0, \theta_0)\in H^2(\mathbb{R}^3)^3$,
  with $\dv u_0 =0$, be such that
  %
   \begin{equation} \label{e:keyii1}
   \ell - C_1\, \big(\|\Lambda^{\frac{1}{2}} (u_0, v_0, \theta_0)  \| -
    \mathcal R(u_0,\alpha) - \tilde{\mathcal R}(v_0,\beta) \big) >0,
    \end{equation}
    and
    %
    \begin{equation} \label{e:keyii2} 2\, \ell - C_2\,\big(\varepsilon
      + \|\Lambda^{\frac{1}{2}} (u_0,v_0,\theta_0)\|^2\big) >0,
    \end{equation}
    where $C_1, C_2>0$ are the constants $C$ in
    \eqref{stima-u-v-H1-chiusa} and \eqref{e:key2}, respectively.
  
    Further, let $(u,v,\theta)(t)$ be the local strong solution of
    \eqref{TCM-gen} and let $T^\ast>0$ be  the first blow-up
    time, as defined in Theorem \ref{th:bu criterion}.
  %
    Then, the following relations hold true
  %
    \begin{equation}
      \label{e:monotonicity1}
      \sup_{0\leq t <T^\ast} \|(u,v,\theta)(t)\|_{\dot H^{\frac12}}
      \le \|(u_0,v_0,\theta_0)\|_{\dot H^{\frac12}}, \ \ \ \
      \sup_{0\leq t <T^\ast} \|(u,v,\theta)(t)\|_{\dot H^1} \le \|(u_0,v_0,\theta_0)\|_{\dot H^1},
    \end{equation}
    and moreover
  %
    \begin{equation}
      \label{e:integrability1}
      \int_{0}^{T^\ast} \|\Delta (u,v,\theta)(t)\|^2\, dt < +\infty.
    \end{equation}
  \end{lem}

\begin{proof}[Proof of Lemma~\ref{lem:monotone}]
  %
  From \eqref{e:keyii1}--\eqref{e:keyii2} and the continuous
  dependence in time with respect to the initial data,
  \eqref{stima-u-v-H1-chiusa} and \eqref{e:key2} imply that
  %
  \begin{equation}
    \label{e:mono} \frac{d}{dt}\| \nabla (u,v,\theta)
    (t)\|^2\leq 0\,\,\, \textrm{ and }\,\,\, \frac{d}{dt} \|
    \Lambda^{\frac12} (u,v,\theta)(t)\|^2 \le 0,
  \end{equation}
  for every $0<t<T_1$, for some $T_1>0$. This implies that
  \eqref{e:keyii1} and \eqref{e:keyii2} hold, for every $0<t<T_1$. We
  claim that \eqref{e:mono} must hold true for every $0<t<T^*$, where
  $T^\ast$ is the (supposed) first blow-up time. Assume then, by
  contradiction, that there exists $T_2>0$ with $T_1 < T_2 <T^*$ such
  that
  %
  \begin{equation*}
    \frac{d}{dt}\|\nabla (u,v,\theta)(T_2)\|^2 >0.
  \end{equation*}
  The case in which is $\frac{d}{dt}\|\Lambda^{\frac{1}{2}} (u,v,\theta)(T_2)\|^2 >0$ can be
  treated similarly.

  Now, set
  %
  \begin{equation*}
    \tilde{T}\doteq \inf\left\{\tau >T_1: \, \frac{d}{dt}\|\nabla (u,v,\theta)(\tau)\|^2 >0\right\}.
  \end{equation*}
  Clearly, $\tilde{T}$ is such that for every $0<t<\tilde{T}$,
  \eqref{e:mono} holds true. This implies that there exists a
  neighborhood of $\tilde{T}$, say
  $(\tilde{T}-\delta, \tilde{T}+\delta)$ such that \eqref{e:keyii1}
  and \eqref{e:keyii2} are true. By using \eqref{stima-u-v-H1-chiusa}
  and \eqref{e:key2}, we deduce that also \eqref{e:mono} is true, in
  $(\tilde{T}-\delta, \tilde{T}+\delta)$. This provides a
  contradiction. Hence, we proved \eqref{e:monotonicity1}.
  
  By integrating \eqref{stima-u-v-H1-chiusa}, with respect to time
  variable $t\in [0, T^\ast)$, we get also \eqref{e:integrability1}.
%
\end{proof}

\end{document}
